En ce début du XVII\ieme{} siècle, il semble donc qu'on a
reconstitué les remparts de notre cité. Le nombre de ses
foyers est singulièrement faible. 25 sont en ville et une
vingtaine sont à l'extérieur : ceux qu'on appelle
« les forains ». Quelle dégringolade si l'on songe aux
quatre cents feux du temps jadis !!! Parmi ces familles,
on trouve des noms qui sont encore portés aujourd'hui...
peut-être par leurs descendants : Paulhe, Trouillet, Puech,
Suc, Augé, Lacombe, Galonnier, Combes, Guiraud, Cadars,
Roussel. Jean Guiraud à la Mouline Ferran\footnote{Moulin de Ferran?} y possédait
alors un moulin « drapier » et un moulin à blé. En ce
moment-là, la « Croux novo » se dressait sur le chemin
d'Alban à Réalmont et celle de « Peyre Bénigne »
(Pierre Benoît) sur celui d'Alban à Peyreblanque.

\par

Si l'on veut remarquer que le 15 juillet 1604, un violent
orage avec grêle, dévasta la région d'Alban, on comprendra
facilement que le pays en ce moment était bien digne de
pitié. Personne ne voulait des charges municipales.
En mai 1600, il n'y eut aucun candidat au consulat.
Le curé de l'époque, pour se procurer une cloche pour
église\footnote{Pour l'église? Pour église? D'église?}, fut obligé d'aliéner un pré paroissial à la famille
Passemar de Saint André. Ces vicomtes de Saint André
portent désormais le titre de barons d'Alban. Notre cité
en garde encore un certain souvenir avec la route dite
de « Saint André ». Le susdit pré ne sera restitué à la
paroisse qu'en 1755. Cent cinquante quatre ans pour
payer une cloche... C'est un record de clochard !!!

\par

Encore que les Saint André se parent du titre de
« barons d'Alban », c'est un Louis d'Albert qui revendique
d'être seigneur d'Alban, puisque c'est lui qu'on convoque
en septembre 1602 pour faire exécuter les édits de paix.
Ce personnage conserve titre et fonction pendant des
années. On le retrouve en 1612 pour passer un contrat
avec le chapitre de Sainte Cécile d'Albi et également
en 1614.